\section{ESPECIFICAÇÃO DETALHADA DO SISTEMA}
\label{ap:especificacao}

\subsection{Requisitos funcionais}

As fontes indicadas na última coluna orientaram o levantamento e a priorização dos requisitos; elas não constituem prescrição literal de cada funcionalidade.

\begin{longtable}{p{0.10\textwidth}p{0.55\textwidth}p{0.10\textwidth}p{0.17\textwidth}}
\caption{Requisitos funcionais do SIGE Desk}\label{tab:requisitos-funcionais}\\
\hline
\textbf{ID} & \textbf{Requisito} & \textbf{Prioridade} & \textbf{Fontes de embasamento} \\
\hline
\endfirsthead
\hline
\textbf{ID} & \textbf{Requisito} & \textbf{Prioridade} & \textbf{Fontes de embasamento} \\
\hline
\endhead
RF-01 & O sistema deve permitir autenticação de usuários. & Alta & \cite{servicedesk2023,brasil2018} \\
RF-02 & O sistema deve controlar permissões conforme o perfil de usuário. & Alta & \cite{servicedesk2023,brasil2018} \\
RF-03 & O administrador deve poder cadastrar, editar, ativar e inativar clientes; clientes inativos não devem ficar disponíveis para novas solicitações. & Alta & \cite{kharisma2024,steponaitis2023} \\
RF-04 & O administrador ou atendimento deve poder cadastrar campanhas vinculadas a um cliente. & Alta & \cite{younas2025,agus2019} \\
RF-05 & O sistema deve permitir abrir ticket com cliente, campanha, canal, tipo, urgência informada, prazo desejado, assunto e descrição. A prioridade oficial e os prazos de SLA serão definidos na triagem. & Alta & \cite{servicerequest2020,servicedesk2023} \\
RF-06 & O sistema deve gerar identificador único para cada ticket. & Alta & \cite{servicedesk2023,iso29148} \\
RF-07 & O atendimento deve poder classificar e atribuir um ticket a um responsável. & Alta & \cite{servicedesk2023,steponaitis2023} \\
RF-08 & O responsável deve poder atualizar o status do ticket. & Alta & \cite{servicerequest2020} \\
RF-09 & O sistema deve registrar data, hora, usuário, motivo, campo alterado, valor anterior e novo valor em cada alteração de status, responsável, prioridade, prazo ou aprovação. & Alta & \cite{servicedesk2023,brasil2018} \\
RF-10 & Usuários autorizados devem poder incluir comentários e anexos no ticket. & Alta & \cite{servicedesk2023} \\
RF-11 & O sistema deve permitir registrar aprovação, reprovação ou solicitação de complemento pelo cliente. & Média & \cite{servicerequest2020,steponaitis2023} \\
RF-12 & O responsável deve poder registrar a ação executada e evidências no encerramento. & Alta & \cite{servicerequest2020} \\
RF-13 & O sistema deve permitir cancelar ticket com motivo obrigatório. & Média & \cite{servicerequest2020} \\
RF-14 & O usuário deve poder filtrar tickets por cliente, campanha, tipo, prioridade, responsável, status e período. & Alta & \cite{abbonizio2023,younas2025} \\
RF-15 & O painel deve apresentar total de tickets por status, prioridade, responsável e prazo. & Média & \cite{abbonizio2023,steponaitis2023} \\
RF-16 & O sistema deve identificar tickets vencidos e tickets aguardando resposta do cliente. & Média & \cite{servicedesk2023,servicerequest2020} \\
RF-17 & O sistema deve permitir consultar o histórico completo de cada ticket. & Alta & \cite{servicedesk2023,brasil2018} \\
RF-18 & O sistema deve permitir registrar métricas de contexto, como CTR, CPC, conversão, CPA e ROAS, quando aplicável. & Baixa & \cite{yahia2024,nain2025} \\
RF-19 & O sistema deve permitir exportar uma lista de tickets filtrada para apoio a relatórios. & Baixa & \cite{abbonizio2023,steponaitis2023} \\
RF-20 & O sistema deve notificar no sistema e, por e-mail quando configurado, os envolvidos autorizados em atribuição, comentário, aguardo, complemento confirmado, validação, vencimento, conclusão, reabertura ou cancelamento. & Alta & \cite{servicedesk2023,servicerequest2020} \\
RF-21 & O sistema deve exibir os prazos de primeira resposta e resolução do SLA, incluindo situação de vencimento ou pausa. & Alta & \cite{servicedesk2023} \\
RF-22 & O administrador deve poder manter tipos de demanda, definindo campos obrigatórios e necessidade de aprovação e evidência. & Alta & \cite{servicerequest2020,iso29148} \\
RF-23 & O administrador deve poder configurar calendário, horário, fuso, feriados, recessos e prazos do SLA por prioridade, com regra opcional por cliente ou tipo. & Alta & \cite{servicedesk2023} \\
\hline
\multicolumn{4}{l}{\footnotesize\textbf{Fonte:} Elaborado pelo grupo a partir do levantamento de requisitos e das fontes indicadas.}
\end{longtable}

\subsection{Requisitos não funcionais}

\begin{longtable}{p{0.10\textwidth}p{0.15\textwidth}p{0.50\textwidth}p{0.20\textwidth}}
\caption{Requisitos não funcionais do SIGE Desk}\label{tab:requisitos-nao-funcionais}\\
\hline
\textbf{ID} & \textbf{Aspecto} & \textbf{Critério inicial} & \textbf{Fontes de embasamento} \\
\hline
\endfirsthead
\hline
\textbf{ID} & \textbf{Aspecto} & \textbf{Critério inicial} & \textbf{Fontes de embasamento} \\
\hline
\endhead
RNF-02 & Segurança & Senhas devem ser armazenadas de forma protegida; autenticação, autorização por perfil e auditoria devem impedir e registrar acessos indevidos. & \cite{brasil2018,iso29148} \\
RNF-03 & Privacidade & O sistema deve coletar somente dados necessários e não deve armazenar bases de audiência ou dados sensíveis. & \cite{brasil2018,daoud2023} \\
RNF-04 & Rastreabilidade & Alterações de status, responsável, prazo e aprovação devem permanecer no histórico. & \cite{servicedesk2023,brasil2018} \\
RNF-05 & Integridade & Um ticket concluído ou cancelado não poderá ser apagado sem manter registro administrativo. & \cite{iso29148,brasil2018} \\
RNF-07 & Compatibilidade & A interface deve funcionar em navegadores modernos. & \cite{iso29148} \\
RNF-09 & Acessibilidade & Campos devem possuir rótulos claros, contraste adequado e navegação possível por teclado. & \cite{radhitya2024} \\
\multicolumn{4}{p{0.95\textwidth}}{\footnotesize\textit{Nota:} Por se tratar de protótipo acadêmico com dados fictícios ou autorizados, retenção, backup, restauração e disponibilidade contínua não integram a entrega. Esses controles deverão ser definidos em eventual adoção real, conforme a LGPD.} \\
\hline
\multicolumn{4}{l}{\footnotesize\textbf{Fonte:} Elaborado pelo grupo a partir do levantamento de requisitos e das fontes indicadas.}
\end{longtable}

\subsection{Regras de negócio}

\begin{longtable}{p{0.12\textwidth}p{0.80\textwidth}}
\caption{Regras de negócio do SIGE Desk}\label{tab:regras-negocio}\\
\hline
\textbf{ID} & \textbf{Regra} \\
\hline
\endfirsthead
\hline
\textbf{ID} & \textbf{Regra} \\
\hline
\endhead
RN-01 & A abertura exige cliente, tipo, urgência informada, assunto, descrição e solicitante. Registro em nome do Cliente guarda separadamente autor e solicitante. Prioridade é obrigatória para executar. \\
RN-02 & Campanha existente deve ser compatível com o cliente. Campanha não cadastrada pode ser identificada na abertura, mas precisa ser vinculada a campanha ativa antes da execução. \\
RN-03 & Somente Atendimento ou Administrador pode definir/alterar prioridade, prazo e responsável; esses campos são confirmados para encaminhar ou retomar execução. \\
RN-04 & A matriz de transições define o ator de cada mudança: Atendimento/Admin tratam abertura, triagem, complemento e reabertura; Gestor atribuído trata execução; Cliente ativo aprova/corrige em Em validação e reabre em Concluída. Atendimento/Admin não aprovam em seu nome. \\
RN-05 & Após a execução, o ticket registra ação e material/evidência quando configurados. Tipo com aprovação segue para Em validação; tipo sem aprovação pode concluir com dispensa derivada da configuração. \\
RN-06 & Concluída exige aprovação quando obrigatória ou dispensa automática do tipo; aprovação obrigatória não pode ser dispensada manualmente. \\
RN-07 & Cancelada exige motivo, confirmação e só pode ser alcançada antes da execução. \\
RN-08 & Aguardando cliente registra informação necessária, origem e início de execução. Complemento recebido só retoma a fase de origem após confirmação de suficiência. \\
RN-09 & Cliente vê apenas sua organização; Atendimento/Admin veem a agência; Gestor vê apenas tickets atribuídos. Menus, anexos e notificações seguem o mesmo escopo. \\
RN-10 & Tipos podem exigir evidência e validação antes da conclusão. Para criativo, o material submetido é evidência anterior à decisão do Cliente. \\
RN-11 & Reabertura de concluída exige justificativa, preserva histórico, inicia novo ciclo e requer confirmação de prioridade/responsável antes da execução. Correção em validação mantém o ciclo. \\
RN-12 & Vencimento sinaliza se é de primeira resposta ou resolução e notifica responsável, Atendimento e Administrador; antes da classificação, prazo fica pendente. \\
RN-13 & Atribuição, comentário, solicitação/confirmação de complemento, validação, conclusão, reabertura e cancelamento notificam envolvidos autorizados no sistema. \\
RN-14 & Confirmação automática não conta como primeira resposta; é necessário retorno efetivo de Atendimento, Administrador ou responsável registrado no ticket. \\
\hline
\multicolumn{2}{l}{\footnotesize\textbf{Fonte:} Elaborado pelo grupo a partir das regras preliminares.}
\end{longtable}

\subsection{Dados principais}

\begin{longtable}{p{0.20\textwidth}p{0.72\textwidth}}
\caption{Entidades e dados iniciais}\label{tab:dados-principais}\\
\hline
\textbf{Entidade} & \textbf{Dados iniciais} \\
\hline
\endfirsthead
\hline
\textbf{Entidade} & \textbf{Dados iniciais} \\
\hline
\endhead
Usuário & Identificador, nome, e-mail, perfil, status e cliente vinculado quando o perfil for Cliente. \\
Cliente & Identificador, nome, contato e status. \\
Campanha & Identificador, cliente, nome, canal, objetivo e status. \\
Tipo de demanda & Identificador, nome, status, campos com rótulo, tipo, obrigatoriedade, opções e validação, e necessidade de aprovação e evidência. \\
Configuração de SLA & Calendário, horário, fuso, feriados, recessos, prioridade, prazos, pausa em validação e escopo padrão, por cliente, tipo ou ambos. \\
Ticket & Número, cliente, campanha ou identificação pendente, tipo/regra aplicada, canal, urgência, prazo desejado, prioridade, SLA, assunto, descrição, solicitante, autor, responsável, status, origem do aguardo e ciclo. \\
Comentário & Ticket, autor, data/hora, texto, anexos e indicação de retorno efetivo, quando aplicável. \\
Histórico & Ticket, campo alterado, valor anterior, novo valor, usuário, data e hora. \\
Aprovação & Ticket, material avaliado, decisão, usuário, data e hora e observação. \\
Evidência & Ticket, descrição, link ou arquivo, autor, data/hora e natureza da evidência. \\
Anexo & Identificador, ticket, nome original, tipo, tamanho, referência de armazenamento, autor, data e hora. \\
Notificação & Ticket, evento, destinatário, canal, resumo, geração, situação de envio e data de leitura, quando aplicável. \\
\hline
\multicolumn{2}{l}{\footnotesize\textbf{Fonte:} Elaborado pelo grupo.}
\end{longtable}

\subsection{Catálogo inicial de tipos de demanda}

\begin{longtable}{p{0.14\textwidth}p{0.24\textwidth}p{0.14\textwidth}p{0.19\textwidth}p{0.14\textwidth}}
\caption{Catálogo inicial de tipos de demanda}\label{tab:tipos-demanda}\\
\hline
\textbf{Tipo} & \textbf{Dados adicionais} & \textbf{Aprovação} & \textbf{Evidência} & \textbf{SLA sugerido} \\
\hline
\endfirsthead
\hline
\textbf{Tipo} & \textbf{Dados adicionais} & \textbf{Aprovação} & \textbf{Evidência} & \textbf{SLA sugerido} \\
\hline
\endhead
Ajuste de orçamento & Campanha ativa ou identificação pendente, orçamento atual/proposto, justificativa e data. & Cliente quando configurado como sensível. & Valor aplicado e captura ou link disponível. & Alta ou Urgente. \\
Alteração ou pausa de anúncio & Campanha ativa ou identificação pendente, canal, ativo, ação e motivo. & Cliente quando houver impacto relevante. & Comentário, captura ou link da alteração. & Alta ou Média. \\
Aprovação de criativo & Campanha ativa ou identificação pendente, canal, criativo e data limite. & Obrigatória. & Material submetido e descrição da versão; decisão somente após validação. & Média. \\
Relatório de desempenho & Campanha ativa ou identificação pendente, período, métricas e formato. & Não obrigatória, salvo configuração. & Arquivo ou link do relatório. & Média. \\
Análise de métricas & Campanha ativa ou identificação pendente, período, questão e indicadores disponíveis. & Não obrigatória. & Comentário analítico e relatório de apoio. & Baixa ou Média. \\
\hline
\multicolumn{5}{l}{\footnotesize\textbf{Fonte:} Elaborado pelo grupo.}
\end{longtable}

\subsection{Matriz de transições de status}

\begin{longtable}{p{0.16\textwidth}p{0.20\textwidth}p{0.22\textwidth}p{0.27\textwidth}}
\caption{Transições permitidas no ciclo de vida do ticket}\label{tab:transicoes-status}\\
\hline
\textbf{Status atual} & \textbf{Próximo status} & \textbf{Quem altera} & \textbf{Registro e efeito no SLA} \\
\hline
\endfirsthead
\hline
\textbf{Status atual} & \textbf{Próximo status} & \textbf{Quem altera} & \textbf{Registro e efeito no SLA} \\
\hline
\endhead
Aberta & Em triagem; Cancelada & Atendimento; Administrador & Registra início da triagem ou motivo/confirmação de cancelamento; prazo fica pendente até classificação. \\
Em triagem & Em execução; Aguardando cliente; Cancelada & Atendimento; Administrador & Para executar, registra prioridade, SLA e responsável; para aguardar, informação e origem; para cancelar, motivo/confirmação. \\
Em execução & Aguardando cliente; Em validação; Concluída & Responsável atribuído & Registra ação e material/evidência; conclusão direta requer dispensa automática do tipo. \\
Aguardando cliente (triagem) & Em triagem; Cancelada & Atendimento; Administrador & Confere complemento e retoma, solicita novo complemento ou cancela antes da execução. \\
Aguardando cliente (execução) & Em execução & Atendimento; Administrador & Confere complemento e retoma execução ou solicita novo complemento; não permite cancelamento. \\
Em validação & Concluída; Reaberta & Cliente ativo da organização & Registra aprovação/correção e material avaliado; correção mantém o ciclo. \\
Concluída & Reaberta & Cliente ativo; Atendimento; Administrador & Exige justificativa e inicia novo ciclo de resolução. \\
Reaberta & Em execução & Atendimento; Administrador & Confirma responsável/prioridade e registra origem da correção ou reabertura. \\
\hline
\multicolumn{4}{l}{\footnotesize\textbf{Fonte:} Elaborado pelo grupo.}
\end{longtable}

\subsection{Matriz de rastreabilidade}

\scriptsize
\begin{longtable}{p{0.06\textwidth}p{0.16\textwidth}p{0.18\textwidth}p{0.23\textwidth}p{0.08\textwidth}p{0.13\textwidth}p{0.06\textwidth}}
\caption{Rastreabilidade preliminar entre necessidades, requisitos e testes}\label{tab:rastreabilidade}\\
\hline
\textbf{ID} & \textbf{Origem} & \textbf{Requisitos} & \textbf{Critério ou teste} & \textbf{Prior.} & \textbf{Responsável} & \textbf{Versão} \\
\hline
\endfirsthead
\hline
\textbf{ID} & \textbf{Origem} & \textbf{Requisitos} & \textbf{Critério ou teste} & \textbf{Prior.} & \textbf{Responsável} & \textbf{Versão} \\
\hline
\endhead
RT-01 & Pedidos dispersos e perda de contexto. & RF-05, RF-06, RF-17. & Ticket identificado e histórico consultável; CT-01 e CT-13. & Alta & Grupo TI SIGE 1 & 1.0 \\
RT-02 & Falta de triagem, prazo e responsável. & RF-07, RF-09, RF-21, RF-23, RN-14. & Triagem e SLA registrados; CT-02 e CT-10. & Alta & Grupo TI SIGE 1 & 1.0 \\
RT-03 & Retrabalho e ausência de aprovação. & RF-11, RF-12, RF-22, RN-05, RN-06, RN-11. & Validação, evidência e reabertura; CT-05 e CT-09. & Alta & Grupo TI SIGE 1 & 1.0 \\
RT-04 & Atrasos e comunicação insuficiente. & RF-16, RF-20, RN-12, RN-13. & Vencimento e eventos notificados; CT-06 e CT-07. & Alta & Grupo TI SIGE 1 & 1.0 \\
RT-05 & Proteção de dados. & RF-02, RNF-02. & Restrição por organização; CT-08. & Alta & Grupo TI SIGE 1 & 1.0 \\
RT-06 & Interface compreensível. & RNF-09. & Campos claros, contraste, navegação por teclado e avaliação de usabilidade. & Alta & Grupo TI SIGE 1 & 1.0 \\
\hline
\multicolumn{7}{l}{\footnotesize\textbf{Fonte:} Elaborado pelo grupo.}
\end{longtable}
\normalsize

\subsection{Critérios de aceite iniciais}

\scriptsize
\begin{longtable}{p{0.23\textwidth}p{0.35\textwidth}p{0.18\textwidth}p{0.10\textwidth}}
\caption{Critérios de aceite e testes relacionados}\label{tab:criterios-aceite}\\
\hline
\textbf{Cenário} & \textbf{Resultado esperado} & \textbf{Requisitos} & \textbf{Teste} \\
\hline
\endfirsthead
\hline
\textbf{Cenário} & \textbf{Resultado esperado} & \textbf{Requisitos} & \textbf{Teste} \\
\hline
\endhead
Cliente abre uma demanda. & Ticket recebe número, status Aberta, urgência, prazo desejado e SLA pendente de classificação; campanha pendente segue para triagem. & RF-05, RF-06, RF-21 & CT-01 \\
Usuário acessa o sistema. & Usuário autenticado acessa somente funções e tickets permitidos. & RF-01, RF-02 & CT-11 \\
Administrador mantém clientes e campanhas. & Cliente pode ser ativado ou inativado e campanha fica vinculada ao cliente. & RF-03, RF-04 & CT-12 \\
Atendimento faz a triagem. & Prioridade, regra/SLA, responsável e primeiro retorno efetivo ficam registrados no histórico. & RF-07, RF-09, RN-03, RN-14 & CT-02 \\
Gestor executa alteração. & Ticket recebe comentário, ação e material/evidência e segue para validação ou conclusão conforme o tipo. & RF-08, RF-10, RF-12, RN-05 & CT-03, CT-05 \\
Cliente envia complemento. & Aguardo registra motivo/origem e só retoma a fase correta após complemento suficiente confirmado. & RF-11, RN-08 & CT-04 \\
Entrega precisa ser confirmada. & Tipo com aprovação segue para Em validação e somente Cliente aprova; tipo sem aprovação conclui com dispensa automática. & RF-11, RN-04, RN-05, RN-06 & CT-05 \\
Cliente solicita correção ou reabre concluída. & Correção mantém o ciclo; reabertura exige justificativa, inicia novo ciclo e requer confirmação antes da execução. & RF-11, RN-11 & CT-05 \\
Prazo é ultrapassado. & Ticket é identificado como vencido. & RF-16, RF-21, RN-12 & CT-06 \\
Evento notificável ocorre. & Envolvidos autorizados recebem notificação no sistema, inclusive complemento confirmado e cancelamento. & RF-20, RN-13 & CT-07 \\
Administrador configura tipo. & Campos e regras de aprovação e evidência são aplicados. & RF-22, RN-05, RN-10 & CT-09 \\
Administrador configura SLA. & Calendário, prazos, pausa e precedência são aplicados a novos tickets; existentes preservam a regra registrada. & RF-23 & CT-10 \\
Usuário consulta histórico. & Consulta exibe autor, data, motivo e valores alterados. & RF-09, RF-17, RNF-04 & CT-13 \\
Cliente tenta acesso indevido. & Acesso é negado e tentativa é registrada. & RF-02, RN-09, RNF-02 & CT-08 \\
\hline
\multicolumn{4}{l}{\footnotesize\textbf{Fonte:} Elaborado pelo grupo.}
\end{longtable}
\normalsize

\subsection{Matriz de notificações}

\begin{longtable}{p{0.20\textwidth}p{0.29\textwidth}p{0.15\textwidth}p{0.24\textwidth}}
\caption{Eventos e destinatários das notificações}\label{tab:notificacoes}\\
\hline
\textbf{Evento} & \textbf{Destinatários} & \textbf{Canal} & \textbf{Conteúdo mínimo} \\
\hline
\endfirsthead
\hline
\textbf{Evento} & \textbf{Destinatários} & \textbf{Canal} & \textbf{Conteúdo mínimo} \\
\hline
\endhead
Ticket atribuído & Responsável e atendimento. & Sistema; e-mail opcional fora do protótipo. & Número, assunto, prioridade e prazo. \\
Comentário incluído & Solicitante, responsável e atendimento autorizados. & Sistema; e-mail opcional fora do protótipo. & Número, autor e resumo. \\
Aguardando cliente & Solicitante e atendimento. & Sistema; e-mail opcional fora do protótipo. & Informação necessária, origem e pausa. \\
Complemento confirmado & Solicitante, responsável e atendimento autorizados. & Sistema; e-mail opcional fora do protótipo. & Informação recebida, fase retomada e prazo. \\
Em validação & Solicitante e atendimento. & Sistema; e-mail opcional fora do protótipo. & Ação, material/evidência e decisão esperada. \\
SLA vencido & Responsável, atendimento e administrador. & Sistema; e-mail opcional fora do protótipo. & Prioridade, tipo de prazo e atraso. \\
Concluída ou Reaberta & Solicitante, responsável e atendimento. & Sistema; e-mail opcional fora do protótipo. & Novo status, justificativa, ciclo e próximo passo. \\
Cancelada & Solicitante, responsável e atendimento. & Sistema; e-mail opcional fora do protótipo. & Motivo, autor e confirmação. \\
\hline
\multicolumn{4}{l}{\footnotesize\textbf{Fonte:} Elaborado pelo grupo.}
\end{longtable}

As notificações não devem revelar conteúdo ou anexos a usuários sem permissão para consultar o ticket.
